\section{问题重述与分析}

\subsection{问题的物理背景}

竞赛平台在半径 $R_{\mathrm{env}}=1800$ m 的圆形区域内随机布置若干干扰源~\cite{cumcm2026b}。机器狗从圆域原点
出发，以 $5$ m/s 的速度移动，通过测向与清除两类动作完成任务：对某频道测向耗时 $5$ s，返回
示向度（方向），且测量误差不超过 $1^\circ$；对某频道实施清除时，只有当机器狗与干扰源的
距离不超过 $20$ m 才能成功。机器狗能听到某干扰源信号的充要条件是"到源的距离不超过该源的
接收半径"，而接收半径 $\rho_{\mathrm{rec}}$ 是每个源固有但未知的量，取值范围为
$[1000,1500]$ m。

问题的第一轮只做了一次测向：机器狗在检测点 $S_1$ 处对某频道测向，得到一个示向度
$\theta_1$。这次测量给出了方向信息，但\textbf{没有给出距离信息}——源可能在示向度射线上
任意远近处。问题二要求在此基础上回答：

\begin{enumerate}
  \item 第二个检测点 $S_2$ 应当按什么策略选择？
  \item 可行而优异的 $S_2$ 构成怎样的候选区域，并说明该区域周边各位置的"适合程度"。
\end{enumerate}

不失一般性，本文取 $S_1$ 为坐标原点、$\theta_1$ 沿 $x$ 轴正方向（即 $\theta_1=0^\circ$）；
所有结果通过对 $S_1$ 的平移与旋转即可搬到任意实际情形。

\subsection{难点的结构化分析}

与"再测一次就完了"的直觉不同，本问有三层困难，且必须同时处理。

\textbf{(1) 源的位置是一个二维不确定集，不是一条线。} 一次测向只排除方向不合的位置，源到
$S_1$ 的距离 $d$ 仍然未知，但有两个硬边界：一是机器狗在 $S_1$ 处确实听到了信号，故
$d\le\rho_{\mathrm{rec}}\le1500$ m；二是题目附件规定距离不超过 $5$ m 时信号过强、测不出
示向度，既然 $S_1$ 测到了示向度，就必有 $d>5$ m。于是源的所有可能位置构成以 $S_1$ 为顶点、
深 $1500$ m、张角 $2^\circ$ 的窄楔形（再与 $1800$ m 圆域求交）。

\textbf{(2) 第二个检测点必须"一定听得到"。} 接收半径未知，若 $S_2$ 选得过远，一旦源的接收
半径偏小，第二次测向将返回无信号，这一趟就白跑了。由于 $\rho_{\mathrm{rec}}$ 的下界是
$1000$ m，要\textbf{保证}第二次测向一定成功，就必须让源不确定集内\textbf{每一个}可能位置
到 $S_2$ 的距离都不超过 $1000$ m。这一条把 $S_2$ 的可选范围压缩成一个面积仅
$0.447$ km$^2$ 的透镜形区域（约为 $1800$ m 圆域的 $0.4\%$）。

\textbf{(3) 定位误差只能用区域口径描述，且必须取最坏情况。} 本项目的定位区域口径与问题一
一致：两条示向度各带 $\pm1^\circ$ 误差，其边界射线相交围出一个四边形（楔形交）。定位精度
用该区域的\textbf{直径}度量：直径小于 $20$ m 才能保证清除成功。由于源在楔形内何处未知、
第二次测向向哪边偏 $1^\circ$ 也未知，选点必须让\textbf{所有这些可能情况中最坏的那一个}
尽可能小——注意"最坏"是两层最坏的复合：源位置取最坏、测向误差取最坏。只做一次测向时，
源方向上的不确定使最坏直径为 $1800$ m（$d$ 从 $5$ m 跨到 $1500$ m，再叠加角向
$\pm1^\circ$），因此第二次测向的价值就体现在能把这一最坏直径压到多小。

\subsection{本文的工作}

本文把上述三层困难写成一条完整的建模与求解链条：

\begin{itemize}
  \item 用\textbf{集员（set-membership）}方式刻画全部不确定性：源的不确定集 $C$ 是楔形，
        测向误差是角度区间，不含任何概率分布假设，因此结论是"保证型"的；
  \item 把"一定听得到"写成\textbf{最坏距离约束}，由此得到可行域（透镜）的显式表达；
  \item 以\textbf{最坏情况定位直径} $J(S_2)$ 为目标函数建立极小化模型，并给出适合度
        $F=J^*/J$ 与候选区域 $\{F\ge1/(1+\eta)\}$ 的定义，直接回答"周边区域的适合度"；
  \item 用\textbf{一阶解析式}解释最优解为何"尽量远、并偏离示向度约 $37^\circ$"；
  \item 引入\textbf{文献判据层}：由两站纯方位测向的 Fisher 信息矩阵导出 CRLB 误差椭圆、
        GDOP 与几何稀释式，证明它们与本文解析式同源、但在本模型最坏情形上系统性偏低约三成，
        因而只用于趋势解释与快速预筛；再据此设计全域认证算法，独立复核最优点；
  \item 给出\textbf{灵敏度分析}（接收半径口径、距离先验、阈值、检测点位置）与
        \textbf{自校验}（解析几何对 shapely、离散化加密、镜像对称、确定性复现）。
\end{itemize}

结论是明确的：第二个检测点应当选在\textbf{距 $S_1$ 约 $1004$ m、相对示向度偏向约
$+37^\circ$} 处（关于示向度镜像的两点等价），最坏定位直径由 $1800$ m 降至 $134.08$ m，
缩小 $13.4$ 倍；符合"最坏直径不超过最优值 $1.1$ 倍"的候选区域是两条窄弧带，合计面积
$0.0207$ km$^2$，仅占可行域的 $4.63\%$。
